% s-158-p158-m002; original PDF 158, 159
\Aphorism{s-158-p158-m002}{243}
Es giebt keine Thatsache, die ich nicht ignoriren, kein Interesse, zu dem ich mich nicht uninteressirt stellen dürfte — vorausgesetzt, dass ich es kann und will. Dieser neutrale Standpunkt, der allein für sich schon Religion und Moral theoretisch widerlegt, ist bisher mit Erbitterung geleugnet worden; jeder „Erzieher der Menschheit“ hat uns gerade das \Foreign{tertium non datur} aufs Gewissen nageln und ein Entweder-Oder abfordern wollen, auf Gebieten, wo uns durchaus Nichts zu einer Parteinahme zwingen kann. Weder die Wissenschaft, noch die allgemeine Wohlfahrt, noch das Heil der Seele oder die Idee der Gattung sind fähig, einen solchen Zwang auszuüben; gesetzt, der Mensch müsse „als Mensch“ sich für Moral interessiren, so bliebe immer noch die Frage zurück, was ihn nöthige, seine Handlungs- oder Gesinnungsmotive gerade in seiner „Menschlichkeit“ zu suchen. Selbst die imperativische Macht des \Sperr{Lebens} fällt für das absolute, Stirner’sche Ich dahin; es zwingt uns nichts, leben zu wollen und aus diesem Willen Consequenzen zu ziehen. Oder folgt etwa aus dem Begriff Egoismus mit logischer Nothwendigkeit der Wille, sein Ego zu erhalten? Vielleicht, aber wer zwingt uns zur \Sperr{Logik}? — In summa: wenn wir uns innerhalb irgend einer Alternative entscheiden, so thun wir es nie, weil wir müssen, sondern weil es uns so beliebt, weil wir freiwillig auf unser „Drittes“, unser tertium und neutrum, verzichten und uns aus Willkür der ebenso willkürlichen Voraussetzung unterordnen, kraft deren die Alternative als solche besteht. Hinsichtlich der religiösen Alternativen hat diese unsere freiwillige Parteinahme nunmehr aufgehört, hinsichtlich anderer dauert sie noch fort: sie von uns erzwingen zu wollen, ist verlorene Mühe!
\par

